\documentclass[11pt]{article}
\usepackage{amsmath, amssymb, amsthm}
\usepackage[margin=1in]{geometry}
\usepackage{enumitem}

\title{Practice Problem Sets: Line Integrals}
\author{Math Curriculum Generator}
\date{}

\begin{document}

\maketitle

\section*{Line Integrals: 30 Practice Problems}

\textbf{Problem 1:} Evaluate the scalar line integral $\int_C \frac{y}{x} \,ds$, where $C$ is the curve parameterized by $x = t^4$, $y = t^3$, for $1 \le t \le 2$.
\\
\textbf{Solution:} First, find $dx/dt = 4t^3$ and $dy/dt = 3t^2$.
Compute $ds = \sqrt{(4t^3)^2 + (3t^2)^2} \,dt = \sqrt{16t^6 + 9t^4} \,dt = t^2\sqrt{16t^2 + 9} \,dt$.
Substitute $x$, $y$, and $ds$ into the integral:
$$ \int_1^2 \frac{t^3}{t^4} \left(t^2\sqrt{16t^2 + 9}\right) \,dt = \int_1^2 t\sqrt{16t^2 + 9} \,dt $$
Use $u$-substitution with $u = 16t^2 + 9$, so $du = 32t \,dt \implies t\,dt = \frac{1}{32}du$.
Limits: $t=1 \implies u=25$; $t=2 \implies u=73$.
$$ \frac{1}{32} \int_{25}^{73} \sqrt{u} \,du = \frac{1}{32} \left[ \frac{2}{3}u^{3/2} \right]_{25}^{73} = \frac{1}{48}\left(73\sqrt{73} - 125\right) $$
\textbf{Final Answer:} $\frac{1}{48}(73\sqrt{73} - 125)$
\vspace{0.5cm}

\textbf{Problem 2:} Evaluate the scalar line integral $\int_C \sqrt{1 + 4x^2} \,ds$, where $C$ is the arc of the parabola $y = x^2$ from $(0,0)$ to $(2,4)$.
\\
\textbf{Solution:} Use $x$ as the parameter. $y = x^2 \implies dy/dx = 2x$.
$ds = \sqrt{1 + (dy/dx)^2} \,dx = \sqrt{1 + 4x^2} \,dx$.
Substitute into the integral:
$$ \int_0^2 \sqrt{1 + 4x^2} \sqrt{1 + 4x^2} \,dx = \int_0^2 (1 + 4x^2) \,dx $$
Evaluate: $\left[ x + \frac{4x^3}{3} \right]_0^2 = 2 + \frac{32}{3} = \frac{38}{3}$.
\\
\textbf{Final Answer:} $\frac{38}{3}$
\vspace{0.5cm}

\textbf{Problem 3:} Evaluate the scalar line integral $\int_C (x+y) \,ds$, where $C$ is defined by $x=2t, y=3t$ for $0 \le t \le 1$.
\\
\textbf{Solution:} $dx/dt = 2$, $dy/dt = 3$. $ds = \sqrt{2^2 + 3^2} \,dt = \sqrt{13} \,dt$.
$$ \int_0^1 (2t + 3t)\sqrt{13} \,dt = \int_0^1 5t\sqrt{13} \,dt = 5\sqrt{13} \left[ \frac{t^2}{2} \right]_0^1 = \frac{5\sqrt{13}}{2} $$
\textbf{Final Answer:} $\frac{5\sqrt{13}}{2}$
\vspace{0.5cm}

\textbf{Problem 4:} Evaluate the scalar line integral $\int_C xyz \,ds$, where $C$ is parameterized by $x = 2t$, $y = t$, $z = 3t$, for $0 \le t \le 1$.
\\
\textbf{Solution:} $ds = \sqrt{2^2 + 1^2 + 3^2} \,dt = \sqrt{14} \,dt$.
Integrand: $(2t)(t)(3t) = 6t^3$.
$$ \int_0^1 6t^3\sqrt{14} \,dt = 6\sqrt{14} \left[ \frac{t^4}{4} \right]_0^1 = \frac{3\sqrt{14}}{2} $$
\textbf{Final Answer:} $\frac{3\sqrt{14}}{2}$
\vspace{0.5cm}

\textbf{Problem 5:} Evaluate the scalar line integral $\int_C \frac{x^2}{y^3} \,ds$, where $C$ is given by $x = t^2$, $y = t$, for $1 \le t \le 2$.
\\
\textbf{Solution:} $dx/dt = 2t$, $dy/dt = 1$. $ds = \sqrt{4t^2 + 1} \,dt$.
Integrand: $\frac{(t^2)^2}{t^3} = \frac{t^4}{t^3} = t$.
$$ \int_1^2 t\sqrt{4t^2 + 1} \,dt $$
Let $u = 4t^2 + 1$, $du = 8t \,dt$. Limits: $t=1 \implies u=5$; $t=2 \implies u=17$.
$$ \frac{1}{8} \int_5^{17} u^{1/2} \,du = \frac{1}{8} \left[ \frac{2}{3}u^{3/2} \right]_5^{17} = \frac{1}{12}(17\sqrt{17} - 5\sqrt{5}) $$
\textbf{Final Answer:} $\frac{1}{12}(17\sqrt{17} - 5\sqrt{5})$
\vspace{0.5cm}

\textbf{Problem 6:} Evaluate $\int_C xy^2 \,ds$, where $C$ is the line segment from $(0,0)$ to $(2,4)$.
\\
\textbf{Solution:} Parametrize the segment: $\mathbf{r}(t) = \langle 2t, 4t \rangle$ for $0 \le t \le 1$.
$ds = \sqrt{2^2 + 4^2} \,dt = \sqrt{20} \,dt = 2\sqrt{5} \,dt$.
$$ \int_0^1 (2t)(16t^2)(2\sqrt{5}) \,dt = 64\sqrt{5} \int_0^1 t^3 \,dt = 64\sqrt{5} \left( \frac{1}{4} \right) = 16\sqrt{5} $$
\textbf{Final Answer:} $16\sqrt{5}$
\vspace{0.5cm}

\textbf{Problem 7:} Evaluate $\int_C xy^2 \,ds$, where $C$ is the right half of the circle $x^2 + y^2 = 16$, oriented counterclockwise.
\\
\textbf{Solution:} Parametrize: $x = 4\cos t$, $y = 4\sin t$, for $-\pi/2 \le t \le \pi/2$.
$ds = \sqrt{(-4\sin t)^2 + (4\cos t)^2} \,dt = \sqrt{16} \,dt = 4 \,dt$.
$$ \int_{-\pi/2}^{\pi/2} (4\cos t)(16\sin^2 t)(4) \,dt = 256 \int_{-\pi/2}^{\pi/2} \sin^2 t \cos t \,dt $$
Let $u = \sin t$, $du = \cos t \,dt$. Limits: $-1$ to $1$.
$$ 256 \int_{-1}^1 u^2 \,du = 256 \left[ \frac{u^3}{3} \right]_{-1}^1 = 256 \left( \frac{1}{3} - \left(-\frac{1}{3}\right) \right) = \frac{512}{3} $$
\textbf{Final Answer:} $\frac{512}{3}$
\vspace{0.5cm}

\textbf{Problem 8:} Evaluate $\int_C (x^2 + y^2) \,ds$, where $C$ is the full circle $x^2 + y^2 = 9$.
\\
\textbf{Solution:} Parametrize: $x = 3\cos t$, $y = 3\sin t$, $0 \le t \le 2\pi$.
$ds = 3 \,dt$. Notice that $x^2 + y^2 = 3^2 = 9$.
$$ \int_0^{2\pi} (9)(3) \,dt = \int_0^{2\pi} 27 \,dt = 54\pi $$
\textbf{Final Answer:} $54\pi$
\vspace{0.5cm}

\textbf{Problem 9:} Evaluate $\int_C z \,ds$, where $C$ is the helix $x = \cos t$, $y = \sin t$, $z = t$, for $0 \le t \le 2\pi$.
\\
\textbf{Solution:} $dx/dt = -\sin t$, $dy/dt = \cos t$, $dz/dt = 1$.
$ds = \sqrt{\sin^2 t + \cos^2 t + 1} \,dt = \sqrt{1 + 1} \,dt = \sqrt{2} \,dt$.
$$ \int_0^{2\pi} t\sqrt{2} \,dt = \sqrt{2} \left[ \frac{t^2}{2} \right]_0^{2\pi} = \sqrt{2} \left( \frac{4\pi^2}{2} \right) = 2\pi^2\sqrt{2} $$
\textbf{Final Answer:} $2\pi^2\sqrt{2}$
\vspace{0.5cm}

\textbf{Problem 10:} Evaluate $\int_C xe^{yz} \,ds$, where $C$ is the line segment from $(0,0,0)$ to $(2,3,4)$.
\\
\textbf{Solution:} $\mathbf{r}(t) = \langle 2t, 3t, 4t \rangle$, $0 \le t \le 1$.
$ds = \sqrt{2^2 + 3^2 + 4^2} \,dt = \sqrt{4 + 9 + 16} \,dt = \sqrt{29} \,dt$.
$$ \int_0^1 (2t)e^{(3t)(4t)}\sqrt{29} \,dt = 2\sqrt{29} \int_0^1 t e^{12t^2} \,dt $$
Let $u = 12t^2$, $du = 24t \,dt$.
$$ 2\sqrt{29} \left( \frac{1}{24} \right) \left[ e^u \right]_0^{12} = \frac{\sqrt{29}}{12}(e^{12} - 1) $$
\textbf{Final Answer:} $\frac{\sqrt{29}}{12}(e^{12} - 1)$
\vspace{0.5cm}

\textbf{Problem 11:} Evaluate the coordinate line integral $\int_C x^2 \,dx + y^2 \,dy$, where $C$ is the curve $y = x^3$ from $x=0$ to $x=2$.
\\
\textbf{Solution:} Let $x=x$. Then $dx = dx$, $y = x^3$, $dy = 3x^2 \,dx$.
$$ \int_0^2 (x^2 + (x^3)^2(3x^2)) \,dx = \int_0^2 (x^2 + 3x^8) \,dx = \left[ \frac{x^3}{3} + \frac{3x^9}{9} \right]_0^2 = \frac{8}{3} + \frac{512}{3} = \frac{520}{3} $$
\textbf{Final Answer:} $\frac{520}{3}$
\vspace{0.5cm}

\textbf{Problem 12:} Evaluate $\int_C (x^3y + \sin x) \,dy$, where $C$ is $y = x^2$ from $x=0$ to $x=\pi$.
\\
\textbf{Solution:} Let $x$ be the parameter. $dy = 2x \,dx$.
$$ \int_0^\pi (x^3(x^2) + \sin x)(2x) \,dx = \int_0^\pi (2x^6 + 2x\sin x) \,dx $$
First integral: $\left[ \frac{2x^7}{7} \right]_0^\pi = \frac{2\pi^7}{7}$.
Second integral requires Integration by Parts: $u=2x$, $dv=\sin x \,dx \implies du=2 \,dx, v=-\cos x$.
$$ [-2x\cos x]_0^\pi - \int_0^\pi -2\cos x \,dx = (-2\pi(-1) - 0) + [2\sin x]_0^\pi = 2\pi $$
\textbf{Final Answer:} $\frac{2\pi^7}{7} + 2\pi$
\vspace{0.5cm}

\textbf{Problem 13:} Evaluate $\int_C e^y \,dx$, where $C$ is the curve $x = y^3$ from $(-1, -1)$ to $(8, 2)$.
\\
\textbf{Solution:} Use $y$ as the parameter. $dx = 3y^2 \,dy$. Limits are $y=-1$ to $y=2$.
$$ \int_{-1}^2 e^y (3y^2 \,dy) $$
Actually, let's reverse to make U-sub clean like the source material: $C: x=y^3$, but evaluate $\int_C e^x \,dy$?
Wait, if it's $\int e^y dx = \int e^y (3y^2 dy)$, this requires integration by parts twice. To mimic the original problem 6 exactly: $\int_C e^x dx$ where $x=y^3 \implies \int e^{y^3} 3y^2 dy$. Let's solve exactly this:
Evaluate $\int_C e^x \,dx$, $C: x=y^3$ from $y=-1$ to $y=2$.
$$ \int_{-1}^2 e^{y^3} (3y^2) \,dy = \left[ e^{y^3} \right]_{-1}^2 = e^8 - e^{-1} $$
\textbf{Final Answer:} $e^8 - \frac{1}{e}$
\vspace{0.5cm}

\textbf{Problem 14:} Evaluate $\int_C z^2 \,dx + x^2 \,dy + y^2 \,dz$, where $C$ is the line segment from $(1,0,0)$ to $(3,1,2)$.
\\
\textbf{Solution:} Parametrize: $x = 1+2t$, $y = t$, $z = 2t$, for $0 \le t \le 1$.
$dx = 2 \,dt$, $dy = 1 \,dt$, $dz = 2 \,dt$.
$$ \int_0^1 ((2t)^2(2) + (1+2t)^2(1) + (t)^2(2)) \,dt = \int_0^1 (8t^2 + 1 + 4t + 4t^2 + 2t^2) \,dt $$
$$ = \int_0^1 (14t^2 + 4t + 1) \,dt = \left[ \frac{14t^3}{3} + 2t^2 + t \right]_0^1 = \frac{14}{3} + 2 + 1 = \frac{23}{3} $$
\textbf{Final Answer:} $\frac{23}{3}$
\vspace{0.5cm}

\textbf{Problem 15:} Evaluate $\int_C y \,dx + x \,dy$, where $C$ is parametrized by $x = \sqrt{t}, y = t$ for $1 \le t \le 4$.
\\
\textbf{Solution:} $dx = \frac{1}{2\sqrt{t}} \,dt$, $dy = 1 \,dt$.
$$ \int_1^4 \left(t \left(\frac{1}{2\sqrt{t}}\right) + \sqrt{t}(1)\right) \,dt = \int_1^4 \left(\frac{1}{2}\sqrt{t} + \sqrt{t}\right) \,dt = \int_1^4 \frac{3}{2} t^{1/2} \,dt $$
$$ = \left[ t^{3/2} \right]_1^4 = 4^{3/2} - 1^{3/2} = 8 - 1 = 7 $$
\textbf{Final Answer:} $7$
\vspace{0.5cm}

\textbf{Problem 16:} Evaluate the vector line integral $\int_C \mathbf{F} \cdot d\mathbf{r}$, where $\mathbf{F} = \langle \sin x, \cos y \rangle$ and $\mathbf{r}(t) = \langle t^2, t^3 \rangle$ for $0 \le t \le 1$.
\\
\textbf{Solution:} $\mathbf{r}'(t) = \langle 2t, 3t^2 \rangle$.
$\mathbf{F}(\mathbf{r}(t)) = \langle \sin(t^2), \cos(t^3) \rangle$.
Dot product: $2t\sin(t^2) + 3t^2\cos(t^3)$.
$$ \int_0^1 (2t\sin(t^2) + 3t^2\cos(t^3)) \,dt = \left[ -\cos(t^2) + \sin(t^3) \right]_0^1 $$
$$ = (-\cos(1) + \sin(1)) - (-\cos(0) + \sin(0)) = 1 - \cos(1) + \sin(1) $$
\textbf{Final Answer:} $1 - \cos(1) + \sin(1)$
\vspace{0.5cm}

\textbf{Problem 17:} Evaluate $\int_C \mathbf{F} \cdot d\mathbf{r}$, where $\mathbf{F} = \langle xy, x+y \rangle$ and $C$ is the line segment from $(0,0)$ to $(2,3)$.
\\
\textbf{Solution:} $\mathbf{r}(t) = \langle 2t, 3t \rangle$ for $0 \le t \le 1$. $\mathbf{r}'(t) = \langle 2, 3 \rangle$.
$\mathbf{F}(\mathbf{r}(t)) = \langle (2t)(3t), 2t+3t \rangle = \langle 6t^2, 5t \rangle$.
Dot product: $12t^2 + 15t$.
$$ \int_0^1 (12t^2 + 15t) \,dt = \left[ 4t^3 + \frac{15t^2}{2} \right]_0^1 = 4 + \frac{15}{2} = \frac{23}{2} $$
\textbf{Final Answer:} $\frac{23}{2}$
\vspace{0.5cm}

\textbf{Problem 18:} Evaluate $\int_C \mathbf{F} \cdot d\mathbf{r}$, where $\mathbf{F} = \langle y, -x \rangle$ and $C$ is the upper half of $x^2 + y^2 = 4$ oriented counterclockwise from $(2,0)$ to $(-2,0)$.
\\
\textbf{Solution:} $\mathbf{r}(t) = \langle 2\cos t, 2\sin t \rangle$, $0 \le t \le \pi$.
$\mathbf{r}'(t) = \langle -2\sin t, 2\cos t \rangle$.
$\mathbf{F}(\mathbf{r}(t)) = \langle 2\sin t, -2\cos t \rangle$.
Dot product: $(2\sin t)(-2\sin t) + (-2\cos t)(2\cos t) = -4\sin^2 t - 4\cos^2 t = -4$.
$$ \int_0^\pi -4 \,dt = -4\pi $$
\textbf{Final Answer:} $-4\pi$
\vspace{0.5cm}

\textbf{Problem 19:} Evaluate $\int_C \mathbf{F} \cdot d\mathbf{r}$, where $\mathbf{F} = \langle x+y, y-z, z^2 \rangle$ and $\mathbf{r}(t) = \langle t, 2t, 3t \rangle$ for $0 \le t \le 1$.
\\
\textbf{Solution:} $\mathbf{r}'(t) = \langle 1, 2, 3 \rangle$.
$\mathbf{F}(\mathbf{r}(t)) = \langle t+2t, 2t-3t, (3t)^2 \rangle = \langle 3t, -t, 9t^2 \rangle$.
Dot product: $(3t)(1) + (-t)(2) + (9t^2)(3) = 3t - 2t + 27t^2 = t + 27t^2$.
$$ \int_0^1 (t + 27t^2) \,dt = \left[ \frac{t^2}{2} + 9t^3 \right]_0^1 = \frac{1}{2} + 9 = \frac{19}{2} $$
\textbf{Final Answer:} $\frac{19}{2}$
\vspace{0.5cm}

\textbf{Problem 20:} Evaluate $\int_C \mathbf{F} \cdot d\mathbf{r}$, where $\mathbf{F} = \langle z, x, y \rangle$ and $\mathbf{r}(t) = \langle t, t^2, t^3 \rangle$ for $0 \le t \le 1$.
\\
\textbf{Solution:} $\mathbf{r}'(t) = \langle 1, 2t, 3t^2 \rangle$.
$\mathbf{F}(\mathbf{r}(t)) = \langle t^3, t, t^2 \rangle$.
Dot product: $t^3(1) + t(2t) + t^2(3t^2) = t^3 + 2t^2 + 3t^4$.
$$ \int_0^1 (3t^4 + t^3 + 2t^2) \,dt = \left[ \frac{3t^5}{5} + \frac{t^4}{4} + \frac{2t^3}{3} \right]_0^1 = \frac{3}{5} + \frac{1}{4} + \frac{2}{3} $$
Common denominator is 60: $\frac{36}{60} + \frac{15}{60} + \frac{40}{60} = \frac{91}{60}$.
\\
\textbf{Final Answer:} $\frac{91}{60}$
\vspace{0.5cm}

\textbf{Problem 21:} Let $C_1$ be the curve $y = x^2$ from $(0,0)$ to $(2,4)$, and let $C_2$ be the exact same curve but oriented from $(2,4)$ to $(0,0)$. Evaluate and compare the scalar line integrals $\int_{C_1} x \,ds$ and $\int_{C_2} x \,ds$.
\\
\textbf{Solution:} Parametrize $C_1$: $x=t, y=t^2$, $0 \le t \le 2$. $ds = \sqrt{1+4t^2} \,dt$.
$\int_{C_1} x \,ds = \int_0^2 t\sqrt{1+4t^2} \,dt = \frac{1}{8}\left[ \frac{2}{3}(1+4t^2)^{3/2} \right]_0^2 = \frac{1}{12}(17\sqrt{17} - 1)$.
Scalar line integrals are independent of path orientation. Therefore, $\int_{C_2} x \,ds$ evaluates to the exact same value. 
\\
\textbf{Final Answer:} Both equal $\frac{1}{12}(17\sqrt{17} - 1)$
\vspace{0.5cm}

\textbf{Problem 22:} Evaluate the coordinate line integral $\int_C y \,dx - x \,dy$ where $C$ is the full circle $x^2 + y^2 = R^2$ oriented counterclockwise.
\\
\textbf{Solution:} $x = R\cos t$, $y = R\sin t$, $0 \le t \le 2\pi$.
$dx = -R\sin t \,dt$, $dy = R\cos t \,dt$.
$$ \int_0^{2\pi} (R\sin t)(-R\sin t) - (R\cos t)(R\cos t) \,dt = \int_0^{2\pi} -R^2(\sin^2 t + \cos^2 t) \,dt $$
$$ \int_0^{2\pi} -R^2 \,dt = -2\pi R^2 $$
\textbf{Final Answer:} $-2\pi R^2$
\vspace{0.5cm}

\textbf{Problem 23:} Evaluate $\int_C (x^2 + y^2 + z^2) \,ds$ where $C$ is $x = t, y = \cos(5t), z = \sin(5t)$ for $0 \le t \le \pi$.
\\
\textbf{Solution:} $ds = \sqrt{1^2 + (-5\sin(5t))^2 + (5\cos(5t))^2} \,dt = \sqrt{1 + 25} \,dt = \sqrt{26} \,dt$.
Integrand: $t^2 + \cos^2(5t) + \sin^2(5t) = t^2 + 1$.
$$ \int_0^\pi (t^2 + 1)\sqrt{26} \,dt = \sqrt{26} \left[ \frac{t^3}{3} + t \right]_0^\pi = \sqrt{26}\left(\frac{\pi^3}{3} + \pi\right) $$
\textbf{Final Answer:} $\sqrt{26}\left(\frac{\pi^3}{3} + \pi\right)$
\vspace{0.5cm}

\textbf{Problem 24:} Evaluate $\int_C ye^{xz} \,ds$ where $C$ is the line segment from $(0,0,0)$ to $(3,2,4)$. 
\\
\textbf{Solution:} $\mathbf{r}(t) = \langle 3t, 2t, 4t \rangle$, $0 \le t \le 1$.
$ds = \sqrt{9 + 4 + 16} \,dt = \sqrt{29} \,dt$.
Integrand: $2t e^{(3t)(4t)} = 2te^{12t^2}$.
This is mathematically identical to Problem 10 due to the commutative property of multiplication.
$$ \int_0^1 2t e^{12t^2} \sqrt{29} \,dt = \frac{\sqrt{29}}{12}(e^{12} - 1) $$
\textbf{Final Answer:} $\frac{\sqrt{29}}{12}(e^{12} - 1)$
\vspace{0.5cm}

\textbf{Problem 25:} Evaluate $\int_C x^2e^x \,dx$ where $C$ is the curve $y = x^2$ from $x=0$ to $x=2$.
\\
\textbf{Solution:} The parameter is $x$. We compute $\int_0^2 x^2 e^x \,dx$.
Integration by parts twice: $u=x^2 \implies 2x \implies 2 \implies 0$. $dv=e^x \implies e^x \implies e^x \implies e^x$.
Antiderivative: $x^2e^x - 2xe^x + 2e^x$.
Evaluate from $0$ to $2$:
$$ (4e^2 - 4e^2 + 2e^2) - (0 - 0 + 2e^0) = 2e^2 - 2 $$
\textbf{Final Answer:} $2e^2 - 2$
\vspace{0.5cm}

\textbf{Problem 26:} \textit{Diagnostic Flaw} - Evaluate $\int_C (x+y) \,ds$ where $C$ is the line segment from $(0,0)$ to $(3,4)$. 
\textbf{Flawed Solution:} 1. $x = 3t, y = 4t$ for $0 \le t \le 1$. 2. Substitute: $\int_0^1 (3t + 4t) \,dt$. 3. Integrate: $\int_0^1 7t \,dt = 7/2$.
\\
\textbf{Error:} The student forgot to compute the arc length differential $ds = \sqrt{(x')^2+(y')^2} \,dt = 5 \,dt$. 
\\
\textbf{Correct Answer:} $\int_0^1 (7t)(5) \,dt = 35 \left[ \frac{t^2}{2} \right]_0^1 = \frac{35}{2}$
\vspace{0.5cm}

\textbf{Problem 27:} \textit{Diagnostic Flaw} - Evaluate the vector line integral $\int_C \mathbf{F} \cdot d\mathbf{r}$ for $\mathbf{F} = \langle y, x \rangle$, where $C$ is the line segment from $(1,1)$ to $(0,0)$.
\textbf{Flawed Solution:} 1. I'll parameterize the segment backward: $\mathbf{r}(t) = \langle 1-t, 1-t \rangle$ for $0 \le t \le 1$. 2. $\mathbf{r}'(t) = \langle -1, -1 \rangle$. 3. Dot product: $(1-t)(-1) + (1-t)(-1) = -2 + 2t$. 4. Integrate: $\int_0^1 (-2+2t) \,dt = -1$. 5. Because the path is reversed, I must flip the sign of the final answer: $-(-1) = 1$.
\\
\textbf{Error:} Parametrizing the curve specifically from $B$ to $A$ using $(1-t)$ inherently builds the reversed orientation into the bounds and derivatives; artificially flipping the sign again at the end makes it incorrect. 
\\
\textbf{Correct Answer:} $-1$
\vspace{0.5cm}

\textbf{Problem 28:} \textit{Diagnostic Flaw} - Evaluate $\int_C \mathbf{F} \cdot d\mathbf{r}$ where $\mathbf{F} = \langle 2x, 2y \rangle$ and $C$ is $\mathbf{r}(t) = \langle \cos t, \sin t \rangle$, $0 \le t \le \pi$.
\textbf{Flawed Solution:} 1. $\mathbf{r}'(t) = \langle -\sin t, \cos t \rangle$. 2. Dot product: $2x(-\sin t) + 2y(\cos t)$. 3. Integral is $\int_0^\pi (-2x\sin t + 2y\cos t) \,dt = [2x\cos t + 2y\sin t]_0^\pi = -2x - 2x = -4x$.
\\
\textbf{Error:} The student failed to substitute the components of $\mathbf{r}(t)$ ($x = \cos t$, $y = \sin t$) into the vector field $\mathbf{F}$ before integrating.
\\
\textbf{Correct Answer:} $\int_0^\pi (-2\cos t\sin t + 2\sin t\cos t) \,dt = \int_0^\pi 0 \,dt = 0$
\vspace{0.5cm}

\textbf{Problem 29:} \textit{Fundamental Theorem of Line Integrals (FTLI)} - Evaluate $\int_C 2xy \,dx + x^2 \,dy$ where $C$ is the wildly complex spiraling path $\mathbf{r}(t) = \langle \ln(t+1), \sin(\frac{\pi t}{2}) \rangle$ from $t=0$ to $t=1$.
\\
\textbf{Solution:} This is a conservative vector field $\mathbf{F} = \langle 2xy, x^2 \rangle$. The potential function is $f(x,y) = x^2y$.
Instead of a nightmare integration, apply FTLI: $\int_C \nabla f \cdot d\mathbf{r} = f(\mathbf{r}(1)) - f(\mathbf{r}(0))$.
Start point ($t=0$): $\langle \ln(1), \sin(0) \rangle = \langle 0, 0 \rangle$.
End point ($t=1$): $\langle \ln(2), \sin(\pi/2) \rangle = \langle \ln 2, 1 \rangle$.
$f(\ln 2, 1) - f(0,0) = (\ln 2)^2(1) - (0)^2(0) = (\ln 2)^2$.
\\
\textbf{Final Answer:} $(\ln 2)^2$
\vspace{0.5cm}

\textbf{Problem 30:} \textit{Fundamental Theorem of Line Integrals (3D)} - Evaluate $\int_C \mathbf{F} \cdot d\mathbf{r}$ where $\mathbf{F} = \langle y^2z^3, 2xyz^3, 3xy^2z^2 \rangle$, and $C$ is given by $\mathbf{r}(t) = \langle \cos(\pi t), t^2, e^{t^2-1} \rangle$ for $0 \le t \le 1$.
\\
\textbf{Solution:} $\mathbf{F}$ is conservative with potential function $f(x,y,z) = xy^2z^3$.
Start point ($t=0$): $\langle \cos(0), 0, e^{-1} \rangle = \langle 1, 0, \frac{1}{e} \rangle$.
End point ($t=1$): $\langle \cos(\pi), 1, e^0 \rangle = \langle -1, 1, 1 \rangle$.
Apply FTLI: $f(-1, 1, 1) - f(1, 0, e^{-1}) = (-1)(1)^2(1)^3 - (1)(0)^2(e^{-1})^3 = -1 - 0 = -1$.
\\
\textbf{Final Answer:} $-1$

\newpage
\section*{Concept Checklist}
The following is a mapping of key conceptual themes, calculation problem types, and tested techniques to the respective problem numbers:

\begin{enumerate}[label={[C\arabic*]}]
    \item \textbf{Scalar Line Integral, direct parametrization (2D):} Problems 1, 3, 5
    \item \textbf{Scalar Line Integral, direct parametrization (3D):} Problems 4, 10
    \item \textbf{Scalar Line Integral, geometric setup (line segment):} Problems 6, 21
    \item \textbf{Scalar Line Integral, geometric setup (circle/arc):} Problems 7, 8
    \item \textbf{Coordinate Line Integral ($P\,dx + Q\,dy$), function substitution:} Problems 11, 12, 13, 15
    \item \textbf{Coordinate Line Integral ($P\,dx + Q\,dy + R\,dz$), direct parametrization:} Problem 14
    \item \textbf{Vector Line Integral ($\int \mathbf{F} \cdot d\mathbf{r}$), 2D:} Problems 16, 17, 18, 22
    \item \textbf{Vector Line Integral ($\int \mathbf{F} \cdot d\mathbf{r}$), 3D:} Problems 19, 20
    \item \textbf{Algebraic Trick - Pythagorean identity under radical for $ds$:} Problems 7, 9, 23
    \item \textbf{Integration Technique - U-substitution for inner functions:} Problems 1, 5, 7, 10, 13, 24
    \item \textbf{Integration Technique - Integration by Parts for Mixed Types:} Problems 12, 25
    \item \textbf{Flaw Check - Forgetting $ds$ magnitude:} Problem 26
    \item \textbf{Flaw Check - Failing to substitute $x,y$ into $\mathbf{F}$:} Problem 28
    \item \textbf{Flaw Check - Reversing orientation properly:} Problems 21, 27
    \item \textbf{Advanced/Untested Concept - FTLI (Conservative Vector Fields):} Problems 29, 30
\end{enumerate}

\end{document}